\capitulo{1}{Introducción}

El aprendizaje automático \cite{bishop2006PatternRecognition} es un subcampo de la inteligencia artificial centrado en el desarrollo de algoritmos y modelos capaces de aprender a partir de datos asociados a problemas específicos. De este modo, pueden realizar predicciones o tomar decisiones sin haber sido programados explícitamente para resolver cada caso concreto.

Entre sus áreas más relevantes se encuentran la clasificación, cuyo objetivo es asignar una etiqueta o categoría a una observación a partir de sus características; por ejemplo, determinar si un correo electrónico es spam o no spam. También destaca la regresión, orientada a predecir valores numéricos continuos, como el precio de una vivienda según sus características y ubicación.

Entre los modelos de aprendizaje automático más empleados, los árboles de decisión han mantenido una gran popularidad durante décadas, debido a su utilidad, flexibilidad y facilidad de interpretación \cite{costa2023RecentAdvancesDecisionTrees}. Estos modelos permiten descomponer decisiones complejas en una sucesión de elecciones simples, representadas de forma gráfica mediante nodos y ramas. Por ello, constituyen una herramienta accesible para profesionales de diferentes áreas que buscan aplicar técnicas de aprendizaje automático sin perder la capacidad de comprender y justificar las decisiones tomadas por el modelo.

No obstante, la construcción de árboles de decisión puede resultar compleja para estudiantes o personas sin experiencia previa en aprendizaje automático. Para comprender correctamente su funcionamiento, es necesario conocer conceptos como la entropía, la ganancia de información, la información de partición, el \textit{Gain Ratio}, la selección de umbrales para atributos numéricos y los procesos de poda. La combinación de estos elementos puede dificultar la comprensión del algoritmo, especialmente cuando se estudia únicamente desde una perspectiva teórica.

El algoritmo C4.5, propuesto por Quinlan como evolución del algoritmo ID3, permite construir árboles de decisión aplicables a problemas de clasificación. Frente a ID3, C4.5 incorpora mejoras relevantes, como el uso del \textit{Gain Ratio} para seleccionar los atributos de partición, el tratamiento de atributos numéricos mediante umbrales, la gestión de valores ausentes y la poda posterior del árbol para reducir el sobreajuste. Estas características hacen que C4.5 sea un algoritmo más completo y aplicable a conjuntos de datos reales, aunque también incrementan la dificultad de su enseñanza y comprensión.

Este algoritmo ha sido utilizado en múltiples estudios y ámbitos de aplicación. Un ejemplo es el artículo \textit{Predicting Students Graduate on Time Using C4.5 Algorithm}~\cite{artIntC4.5}, en el que se analiza la capacidad de C4.5 para predecir si estudiantes universitarios finalizarían sus estudios dentro del plazo previsto o con retraso. A partir de los registros académicos del alumnado correspondientes a los años 2021 y 2022, el estudio obtuvo una efectividad del 92,6\%, lo que refleja la utilidad del algoritmo para abordar problemas reales de clasificación basados en datos.

Los simuladores didácticos han demostrado ser recursos eficaces de apoyo en el proceso de enseñanza-aprendizaje en distintas áreas de la informática. Estas herramientas se emplean, por ejemplo, para simular algoritmos de planificación de CPU en sistemas operativos \cite{gonzalezRodriguez2024CPUScheduling}, explicar algoritmos sobre grafos en estructuras de datos \cite{antolinSendino2003Sigraf} o representar el comportamiento de robots \cite{gonzalezGarcia2023SimuladorRobotsArduino}. Su principal ventaja consiste en permitir que el alumnado experimente de manera práctica con los conceptos teóricos, observando de forma visual las consecuencias de cada decisión tomada durante la ejecución de un algoritmo.

Aunque existen herramientas y recursos educativos relacionados con árboles de decisión, muchas de ellas se centran en mostrar el árbol final generado o en explicar de forma limitada conceptos como la entropía y la ganancia de información. Son menos frecuentes las aplicaciones que permiten estudiar de manera guiada y detallada todas las etapas de construcción de un árbol mediante C4.5, incluyendo el cálculo del \textit{Gain Ratio}, la comparación de posibles umbrales en atributos numéricos y el proceso de poda del árbol resultante.

Por ello, se pretende crear una herramienta interactiva que permita enseñar de forma eficaz el funcionamiento del algoritmo C4.5. Como referencia se toma el proyecto de enseñanza interactiva del algoritmo \texttt{ID3}, denominado \texttt{TFG-decision-trees-sim}\footnote{TFG-decision-trees-sim: \url{https://github.com/danieldf01/TFG-decision-trees-sim}}, el cual constituye una herramienta eficaz para adquirir los conocimientos necesarios para comprender dicho algoritmo. A partir de esta referencia, se busca ampliar el enfoque didáctico hacia C4.5, incorporando las características propias de este algoritmo y facilitando su estudio de manera visual e interactiva.

Una aplicación interactiva que permita visualizar cada paso del proceso de construcción puede ayudar a comprender no solo cuál es el árbol final obtenido, sino también por qué se selecciona un atributo concreto, cómo se divide el conjunto de datos y de qué forma se eliminan ramas innecesarias mediante la poda.

En este trabajo se presenta una herramienta docente diseñada específicamente para mejorar la comprensión de la construcción de árboles de decisión mediante el algoritmo C4.5. La aplicación permite al usuario visualizar e interactuar con las distintas fases del algoritmo, facilitando el aprendizaje de conceptos fundamentales como la entropía, la ganancia de información, la información de partición, el \textit{Gain Ratio}, la selección de umbrales para atributos numéricos y la poda del árbol. Asimismo, permite utilizar conjuntos de datos de ejemplo o introducir datos propios, favoreciendo un aprendizaje práctico y adaptable a distintos contextos educativos.

La herramienta desarrollada pretende contribuir a la enseñanza de los árboles de decisión desde una perspectiva visual, interactiva y didáctica, facilitando la comprensión de uno de los algoritmos más relevantes dentro del aprendizaje automático supervisado.
